% Experiment 13: Implementation of Digital Signature Standard (DSS) for Secure Authentication
\section{Theory}
The Digital Signature Standard (DSS) provides message authenticity and integrity
without confidentiality, relying on the discrete logarithm problem. Global
public parameters are a large prime $p$, a prime divisor $q$ of $(p-1)$, and a
generator $g$ of order $q$ modulo $p$. Each signer has a private key $x$
($0 < x < q$) and public key $y = g^x \bmod p$. To sign a message $m$, compute
hash $H(m)$, choose a random $k$, and generate the signature $(r, s)$:
\[
	r = (g^k \bmod p) \bmod q, \qquad s = k^{-1}(H(m) + x r) \bmod q
	\]
	To verify, the recipient computes $w = s^{-1} \bmod q$,
$u_1 = H(m) w \bmod q$, $u_2 = r w \bmod q$, and
$v = ((g^{u_1} y^{u_2}) \bmod p) \bmod q$. The signature is valid if and
only if $v = r$.

\section{Algorithm}
\begin{enumerate}
\item Select global parameters: prime $p$, prime divisor $q$ of $(p-1)$, and generator $g$.
\item Generate a private key $x$ with $0 < x < q$.
\item Compute the public key $y = g^x \bmod p$.
\item To sign message $m$: compute hash $H(m)$ and choose random $k$ with $\gcd(k,q)=1$.
\item Compute $r = (g^k \bmod p) \bmod q$.
\item Compute $s = k^{-1}(H(m) + xr) \bmod q$. The signature is $(r, s)$.
\item To verify: compute $w = s^{-1} \bmod q$, $u_1 = H(m)w \bmod q$, $u_2 = rw \bmod q$.
\item Compute $v = ((g^{u_1} y^{u_2}) \bmod p) \bmod q$.
\item Accept the signature as valid if $v = r$, else reject it.
\end{enumerate}

\section{Code}
\begin{minted}{cpp}
#include<bits/stdc++.h>
using namespace std;
class dss {
private:
  long long p, q, g, x, y;
  long long modPow(long long base, long long exp, long long mod) {
    long long result = 1;
    base %= mod;
    while (exp > 0) {
      if (exp & 1) result = (__int128)result * base % mod;
      base = (__int128)base * base % mod;
      exp >>= 1;
    }
    return result;
  }
  long long gcd(long long a, long long b) {
    return b == 0 ? a : gcd(b, a % b);
  }
  long long modInverse(long long a, long long m) {
    long long m0 = m, x0 = 0, x1 = 1;
    a = ((a % m) + m) % m;
    while (a > 1) {
      long long qt = a / m;
      long long t = m;
      m = a % m;
      a = t;
      t = x0;
      x0 = x1 - qt * x0;
      x1 = t;
    }
    if (x1 < 0) x1 += m0;
    return x1;
  }
  long long simpleHash(long long m) {
    return (m * 2654435761LL) % q;
  }
public:
  dss(long long primeP, long long primeQ, long long generator, long long privateKey);
  pair<long long, long long> sign(long long msg, long long k);
  bool verify(long long msg, long long r, long long s);
};
dss::dss(long long primeP, long long primeQ, long long generator, long long privateKey) {
  p = primeP;
  q = primeQ;
  g = generator;
  x = privateKey;
  y = modPow(g, x, p);
  cout << "Public key y = g^x mod p = " << y << endl;
}
pair<long long, long long> dss::sign(long long msg, long long k) {
  long long H = simpleHash(msg);
  long long r = modPow(g, k, p) % q;
  long long kInv = modInverse(k, q);
  long long s = (kInv * (H + x * r)) % q;
  cout << "H(m) = " << H << endl;
  cout << "r = (g^k mod p) mod q = " << r << endl;
  cout << "s = k^-1 * (H(m) + x*r) mod q = " << s << endl;
  return make_pair(r, s);
}
bool dss::verify(long long msg, long long r, long long s) {
  long long H = simpleHash(msg);
  long long w = modInverse(s, q);
  long long u1 = (H * w) % q;
  long long u2 = (r * w) % q;
  long long v = (modPow(g, u1, p) * modPow(y, u2, p)) % p % q;
  cout << "w  = s^-1 mod q = " << w << endl;
  cout << "u1 = H(m)*w mod q = " << u1 << endl;
  cout << "u2 = r*w mod q = " << u2 << endl;
  cout << "v  = ((g^u1 * y^u2) mod p) mod q = " << v << endl;
  return v == r;
}
int main() {
  cout << "Enter prime p: ";
  long long p;
  cin >> p;
  cout << "Enter prime divisor q of (p-1): ";
  long long q;
  cin >> q;
  cout << "Enter generator g (of order q mod p): ";
  long long g;
  cin >> g;
  cout << "Enter private key x (0 < x < q): ";
  long long x;
  cin >> x;

  dss signer(p, q, g, x);

  cout << "\nEnter message (as integer, will be hashed): ";
  long long msg;
  cin >> msg;
  cout << "Enter random per-signature value k (0 < k < q, gcd(k,q)=1): ";
  long long k;
  cin >> k;

  cout << "\n--- Signature Generation ---" << endl;
  pair<long long, long long> sig = signer.sign(msg, k);
  cout << "Signature = (r, s) = (" << sig.first << ", " << sig.second << ")" << endl;

  cout << "\n--- Signature Verification ---" << endl;
  bool valid = signer.verify(msg, sig.first, sig.second);

  cout << "\nResult: "
  << (valid ? "SIGNATURE VALID (v == r)" : "SIGNATURE INVALID (v != r)")
  << endl;
  return 0;
}
\end{minted}

\section{Result and Discussion - I/O}
\textbf{Sample Input:}
\begin{verbatim}
Enter prime p: 23
Enter prime divisor q of (p-1): 11
Enter generator g (of order q mod p): 4
Enter private key x (0 < x < q): 9
Enter message (as integer, will be hashed): 3
Enter random per-signature value k (0 < k < q, gcd(k,q)=1): 7
\end{verbatim}

\textbf{Sample Output:}
\begin{verbatim}
Public key y = g^x mod p = 13

--- Signature Generation ---
H(m) = 8
r = (g^k mod p) mod q = 8
s = k^-1 * (H(m) + x*r) mod q = 2
Signature = (r, s) = (8, 2)

--- Signature Verification ---
w  = s^-1 mod q = 6
u1 = H(m)*w mod q = 4
u2 = r*w mod q = 4
v  = ((g^u1 * y^u2) mod p) mod q = 8

Result: SIGNATURE VALID (v == r)
\end{verbatim}

\textbf{Observation:} The computed verification value $v$ equals $r$,
confirming that the digital signature is valid. This demonstrates that a
recipient with only the public key $(p, q, g, y)$ can verify a message
signed with the corresponding private key $x$, without ever learning $x$
itself.

\section{Conclusion}
The DSS was implemented, demonstrating both signing and verification. The
experiment shows that modifying the message or signature causes verification
to fail, confirming authentication and integrity. Production use requires a
secure hash function, larger primes, and a secure random $k$ (since $k$
reuse exposes the private key).
